\section*{Sources and evidence}
\addcontentsline{toc}{section}{Sources and evidence}
The presentation, reference code, worked derivations, generated tables and
illustrations are original teaching artifacts, not claims of new automata
theorems. Classical definitions and the cited experimental architectures
remain attributed below. On 25 September 2026 the Watson--Crick automata
paper's introduction, preliminaries, acceptance and determinism definitions
were inspected. Our empty-input and nonempty-transition conventions are
stated explicitly; no hierarchy theorem is transferred to this variant.
The Benenson paper's mechanism description and Figure 2 caption were inspected
in the author-course archive copy. The mechanism plate is an original
conceptual abstraction, not a reproduction, sequence design or wet-lab protocol.
The Adar paper's abstract and Results discussion of competing transitions
and calibrated programming were inspected. The 2003 Benenson paper was
available here at abstract depth only; its reuse claim is the limited point
used. We do not reconstruct either experiment or import their performance
numbers. The renewal/race and parity-noise examples are declared mathematical
models with assigned parameters. Computed configurations and probabilities
are not measured chemical trajectories. Specialist scientific review and
reader testing remain open; author checks do not replace them.
\begingroup\let\chapter\section
\begin{thebibliography}{9}
\bibitem{czeizler} E. Czeizler, E. Czeizler, L. Kari and K. Salomaa.
\emph{On the descriptional complexity of Watson--Crick automata}.
Theoretical Computer Science 410 (2009), 3250--3260.
\url{https://doi.org/10.1016/j.tcs.2009.05.001}.
Author copy: \url{https://cs.uwaterloo.ca/~lila/pdfs/Watson_Crick_automata.pdf}.
\bibitem{benenson} Y. Benenson, T. Paz-Elizur, R. Adar, E. Keinan,
Z. Livneh and E. Shapiro.
\emph{Programmable and autonomous computing machine made of biomolecules}.
Nature 414 (2001), 430--434.
\url{https://doi.org/10.1038/35106533}.
Mechanism inspected in the Duke molecular-computing course archive copy.
\bibitem{fuel} Y. Benenson, R. Adar, T. Paz-Elizur, Z. Livneh and E. Shapiro.
\emph{DNA molecule provides a computing machine with both data and fuel}.
Proceedings of the National Academy of Sciences 100 (2003), 2191--2196.
\url{https://doi.org/10.1073/pnas.0535624100}.
Abstract: \url{https://pubmed.ncbi.nlm.nih.gov/12601148/}.
\bibitem{adar} R. Adar, Y. Benenson, G. Linshiz, A. Rosner,
N. Tishby and E. Shapiro.
\emph{Stochastic computing with biomolecular automata}.
Proceedings of the National Academy of Sciences 101 (2004), 9960--9965.
\url{https://doi.org/10.1073/pnas.0400731101}.
Open text: \url{https://pmc.ncbi.nlm.nih.gov/articles/PMC454388/}.
\end{thebibliography}\endgroup
